% =====================================================================================================
% vygotskian_experiments.tex — The Vygotskian machine, part 5: experimental evidence
%
% DRAFT (2026-11-20). STRUCTURE ONLY: no experiment of the Vygotskian machine has been run yet. Every number is a
% red \result{} placeholder; the prose says what each experiment measures and which prediction it tests, so that the
% values can be swapped in without rewriting. Order = the build order of the design note (Part E): each experiment
% gates the next. Include after vygotskian_cultural_plane.tex.
% =====================================================================================================

\section{Experimental Evidence}
\label{sec:experiments}

The experiments are presented in the order in which the machine is built, and each experiment must succeed before
the next one begins. Each new piece of the machine is first run in a read-only mode, in which it logs what it would
have done and is scored against what actually happened; only when that score is good is the piece allowed to
influence the robot. The existing metaconcepts were introduced by the same rule (Section~\ref{sec:cortex}), and as a
consequence every experiment below also checks that the previous step changed nothing it should have left untouched.
Table~\ref{tab:exp_overview} lists the experiments together with the predictions that each one tests.

\begin{table}[ht]
    \centering
    \small
    \begin{tabular}{|>{\raggedright\arraybackslash}p{0.06\linewidth}|>{\raggedright\arraybackslash}p{0.36\linewidth}|>{\raggedright\arraybackslash}p{0.14\linewidth}|>{\raggedright\arraybackslash}p{0.32\linewidth}|}
        \hline
        & \textbf{Question} & \textbf{Predictions} & \textbf{Mode} \\ \hline
        E1 & Does culture, elicited from the VLM, agree with the priors we wrote by hand? & L0 & offline \\ \hline
        E2 & Does the function evidence recover what each room is for? & L1 & offline, on recorded tours \\ \hline
        E3 & Does the level-4 function agent work live, and survive a missing object? & L1 & live, read-only \\ \hline
        E4 & Does every mission step predict its belief outcome, and is a divergence caught early? & L3 & live, display only \\ \hline
        E5 & Does the inner voice resolve what the automatic level cannot? & P1, P4, P6 & live \\ \hline
        E6 & Does imagination prepare the robot for what comes next? & L4, P5 & live \\ \hline
        E7 & Is the scaffold withdrawn, and is culture what does the work? & L6, P7, P8 & live, repeated sessions \\ \hline
    \end{tabular}
    \caption{The experiments, in build order. Each gates the next.}
    \label{tab:exp_overview}
\end{table}

\subsection{Setting}
\label{sec:experiments:setting}

All experiments run on \cortex{}, as described in Section~\ref{sec:cortex}, with every agent of
Figure~\ref{fig:cortex_overview} active. The robot operates in a simulated apartment of three rooms connected in a
ring by three doors: an open-plan room containing a kitchen run, a dining table with chairs and a desk; a bedroom; and
a bathroom. Because the simulation fixes the body, the sensors and the world, two versions of the machine can be
compared in the same closed loop on identical tours, and the simulator provides the true function of every room and
the true position of every object. The language agent uses the deployed open-weights model (Section~\ref{sec:cortex}).
The same VLM is used in every condition, so that any differences between conditions arise from how the model is used
rather than from which model is used. During every run, a passive recorder logs the history of the shared graph, so
that the offline experiments can replay real tours. Experiments E5--E7 are repeated on the physical robot, Shadow, in a
real apartment, in order to check that the results obtained in simulation carry over.

\subsection{E1 --- Culture elicited from the VLM}
\label{sec:experiments:e1}

In the current system, the kitchen metaconcept weighs each candidate member of a kitchen unit by its log-odds of
membership, $\log P(c \mid \text{kitchen unit}) / P(c \mid \text{not})$, a value that was written by hand for each of
seven object classes. In this experiment, we elicit the same quantity from the VLM. The model is asked one question per
class and per hypothesis, chooses its answer from a closed vocabulary and runs at temperature zero; each question is
posed in three paraphrases, of which we keep the median answer and record the spread. Table~\ref{tab:exp_e1} compares
the elicited values with the authored ones. Prediction L0 holds if the elicited values agree in sign with the authored
ones, and if replaying the logged kitchen fits with the elicited values changes no correct membership decision.

\begin{table}[ht]
    \centering
    \small
    \begin{tabular}{|l|r|r|r|}
        \hline
        \textbf{Class} & \textbf{Authored} & \textbf{Elicited (median)} & \textbf{Spread} \\ \hline
        cabinet          & $0.2$  & \result{--} & \result{--} \\ \hline
        refrigerator     & $2.0$  & \result{--} & \result{--} \\ \hline
        counter          & $2.5$  & \result{--} & \result{--} \\ \hline
        countertop       & $2.5$  & \result{--} & \result{--} \\ \hline
        hood             & $3.0$  & \result{--} & \result{--} \\ \hline
        microwave        & $1.5$  & \result{--} & \result{--} \\ \hline
        chest of drawers & $-0.5$ & \result{--} & \result{--} \\ \hline
    \end{tabular}
    \caption{Kitchen membership log-odds written by hand, compared with those elicited from the VLM (E1). Number of
    membership decisions changed on replay: \result{--} of \result{--}.}
    \label{tab:exp_e1}
\end{table}

When the same procedure is extended to every class known to the concept agents and to the seven room functions, it
produces the pseudo-counts of Eq.~\eqref{eq:lh_beta_prior}. Their spread across paraphrases is reported for each
entry, because a prior that the VLM cannot state consistently should carry less weight.

\subsection{E2 --- Function evidence on recorded tours}
\label{sec:experiments:e2}

In this experiment, the function evidence $\Lambda_f(r)$ of Section~\ref{sec:lang_hierarchy:culture}, that is, the
evidence that room $r$ serves function $f$, is computed offline over recorded tours from the objects present in each
room and from the places where the robot searched. Each room is labelled with its true function from the simulator,
using the robot's true position at the moment the room was first instantiated, since the robot names rooms in the
order in which it enters them. Two results are reported. The first is the set of functions whose evidence rises above
the null value in each room, compared with the true functions: \result{--} of \result{--} rooms are correct, with
\result{--} false functions. The second is the evolution of $\Lambda_f$ as the exploration of the room becomes more
complete. If the term for objects that are expected but absent works as intended, every $\Lambda_f$ starts at zero and
moves away from zero only as the room is explored (Figure~\result{--}).

\subsection{E3 --- The function agent, live and read-only}
\label{sec:experiments:e3}

In this experiment, the level-4 function agent runs live on the metaconcept chassis while the robot explores the
apartment from scratch. The agent publishes nothing and only logs $\Lambda_f$ for every room. It is scored as in E2,
but now on the posteriors as they evolve during exploration. In a second condition, the refrigerator is removed from
the simulated kitchen. Prediction L1 holds if the kitchen hypothesis becomes weaker without disappearing; we report the
final value of $\Lambda_{\text{kitchen}}$ with the refrigerator (\result{--}), without it (\result{--}), and the null
value (\result{0}). If a function hypothesis collapses when a single diagnostic object is missing, it has been given
too strong a prior, and in that case the priors obtained in E1 are revised before the next experiment begins.

\subsection{E4 --- Plans that predict their belief outcome}
\label{sec:experiments:e4}

Each mission step declares the change of belief that it is expected to produce
(Section~\ref{sec:lang_hierarchy:sequencing}), and the interpreter checks this expectation against the graph. In this
experiment, the demand signal, which decides when the robot switches into focused mode, is computed and displayed but
does not drive anything. Divergences between the plan and the world are seeded deliberately during missions: a door on
which the plan relies is closed, an object is moved, or a room is left unexplored. For each seeded divergence, we
measure whether it is detected, how much earlier it is detected than the moment at which a step would have failed, and
whether the demand signal crossed into focused mode. Prediction L3 states that, when expected outcomes and assumptions are included in the
demand signal, \result{--}\% of the seeded divergences are caught before any step fails, compared with \result{--}\%
when they are not included.

\subsection{E5 --- The inner voice}
\label{sec:experiments:e5}

In this experiment, the inner voice is enabled. Detectors raise questions that the robot addresses to itself about
unexplained occupancy and about unfamiliar objects, the VLM proposes tests that could answer them, and only perception
is allowed to change the answer (Section~\ref{sec:inner_speech:questions}). The machine is compared with \cortex{}
alone on the same tours, in which known objects are moved and unfamiliar ones are introduced. Three measures
correspond to the predictions. The first is the unexplained residual mass per room over time, that is, the part of the
observed occupancy that no current model accounts for (P1: \result{--} with the voice, \result{--} without). The second
is the number of inner questions per hour over repeated visits to the same rooms, while answer accuracy is maintained
(P4: from \result{--} to \result{--}). The third repeats the same loop with the VLM replaced by a fixed set of
hypotheses, namely the existing classes and ``unknown'', in order to measure what the VLM contributes (P6: questions
resolved outside the built-in models, \result{--} versus \result{--}).

\subsection{E6 --- Imagination}
\label{sec:experiments:e6}

While demand remains low, the machine engages in imagination. It rehearses the next legs of the running mission and
completes scenes that it expects but has not yet seen, applying the robot's own models to an imagined copy of the world
(Section~\ref{sec:lang_hierarchy:imagination}). Every imagined prediction is logged and is scored once the
corresponding observation arrives, which measures how well calibrated imagination is for each kind of prediction
(\result{--}). A surprise is then seeded in the middle of a mission, and we measure the latency until a validated new
plan is available, with and without the contingency plans prepared during the periods of low demand (L4, P5:
\result{--}~s versus \result{--}~s), together with the number of views needed to confirm an expected object
(\result{--} versus \result{--}).

\subsection{E7 --- The scaffold withdrawn}
\label{sec:experiments:e7}

The last experiment tests the claim expressed in the name of the machine. The robot lives in the same apartment over
repeated sessions and is given a mix of recurring and novel tasks. If the machine internalises what it receives, three
effects should appear together. First, the weight of the VLM's priors is overtaken by the robot's own counts, and the
calibration $\beta_f$ improves. Second, the number of VLM calls per hour falls while task success is maintained (L6:
from \result{--} to \result{--} calls per hour, success \result{--}). Third, VLM use rises again when tasks are made
harder or new, in the same way as children's private speech does (P8, Figure~\result{--}). Finally, a non-Vygotskian
control condition replaces the culture supplied by the VLM with a library learned only from the robot's own experience
(Section~\ref{sec:inner_speech:library}). Prediction P7 holds if the control matches the machine supplied with culture
on recurring tasks and falls behind it in proportion to how far a task departs from what the robot has seen
(Figure~\result{--}). If such a gap appears, it measures the contribution of culture.

\subsection{What would count against the machine}
\label{sec:experiments:against}

Each experiment can fail in a way that would change the design. If E2 or E3 fails, culture does not belong at level 4
in its present form, and none of the later components is coupled to it. If E4 finds no divergence earlier than a
failed step, the demand signal has no early input, and the inner voice is recruited only after a failure has occurred.
If E6 shows no gain in latency or in the number of views, imagination does not justify its computational cost, and the
robot should simply remain idle when the automatic level suffices. Finally, if in E7 the control without culture keeps
up with the machine supplied with culture on novel tasks, the machine works but is not Vygotskian, because its gains
would then come from having a language rather than from having a culture.
